% =============================================================
% บทที่ 2 - งานวิจัยและทฤษฎีที่เกี่ยวข้อง
% โครงงาน: การออกแบบและพัฒนาเกมเดินเรือด้วยระบบสั่งการด้วยเสียงเพื่อส่งเสริมทักษะการสื่อสาร
% =============================================================

% -----------------------------------------------------------
\section{งานวิจัย โครงงาน และระบบงานที่เกี่ยวข้อง}
% -----------------------------------------------------------

ในส่วนนี้นำเสนองานวิจัยและเกมที่มีลักษณะคล้ายคลึงกับโครงงานนี้ เพื่อศึกษาแนวคิด จุดเด่น และข้อจำกัด แล้วนำมาเป็นแนวทางในการออกแบบและพัฒนาระบบ

\subsection{การเรียนรู้ภาษาผ่านเกมด้วยเทคโนโลยีการประมวลผลภาษาธรรมชาติ}

งานวิจัยของ Sadigzade~\cite{sadigzade2025} นำเสนอการใช้เกมดิจิทัลร่วมกับเทคโนโลยีการรู้จำเสียงและการประมวลผลภาษาธรรมชาติ เพื่อสนับสนุนการเรียนรู้ภาษาอังกฤษ การให้ผู้เรียนได้โต้ตอบด้วยเสียงกับตัวละครปัญญาประดิษฐ์ภายในเกมช่วยสร้างสภาพแวดล้อมจำลองที่สมจริง ผลการศึกษาพบว่าช่วยเพิ่มแรงจูงใจ ลดความวิตกกังวลในการพูด และพัฒนาความคล่องแคล่วในการสื่อสารได้อย่างมีประสิทธิภาพ อย่างไรก็ตาม เกมในลักษณะดังกล่าวมักไม่ได้เน้นความกดดันด้านเวลาหรือสถานการณ์จำลองที่ต้องใช้ไหวพริบอย่างรวดเร็ว

\subsection{เครื่องมือเกมควบคุมด้วยเสียงเพื่อเสริมสร้างการเรียนรู้}

งานวิจัยของ Varma และคณะ~\cite{varma2025} มุ่งเน้นพัฒนาเกมที่ควบคุมการทำงานด้วยคำสั่งเสียงแบบเรียลไทม์ โดยทำงานผ่านการแปลงเสียงพูดเป็นคำสั่งเพื่อควบคุมการเคลื่อนที่ขององค์ประกอบในเกมแทนการใช้อุปกรณ์ควบคุมทั่วไป การควบคุมผ่านเสียงช่วยกระตุ้นให้ผู้เล่นต้องคิดตัดสินใจอย่างรวดเร็วและเปล่งเสียงให้ชัดเจน ซึ่งช่วยพัฒนาทักษะการตอบสนองได้ดี

\subsection{การสำรวจการโต้ตอบด้วยเสียงกับตัวละคร NPC ในเกม}

Zargham และคณะ~\cite{zargham2024} ได้สำรวจความคิดเห็นของผู้เชี่ยวชาญเกี่ยวกับการใช้เสียงพูดเพื่อสื่อสารกับตัวละคร NPC ในวิดีโอเกม ผลการวิจัยระบุว่าการสั่งการด้วยเสียงช่วยเพิ่มความรู้สึกดื่มด่ำ (Immersion) และความสมจริง นอกจากนี้ การให้ NPC ตอบสนองต่อคำสั่งเสียงอย่างเป็นธรรมชาติยังช่วยเพิ่มความผูกพันได้ แต่จำเป็นต้องคำนึงถึงบริบทและความแม่นยำของระบบรู้จำเสียงเพื่อป้องกันความหงุดหงิดจากการแปลความหมายผิดพลาด

\subsection{เกมเชิงพาณิชย์ที่สั่งการด้วยเสียง}

ในอุตสาหกรรมเกมเชิงพาณิชย์ มีการนำระบบรู้จำเสียงมาใช้ในเกมแนวจำลองสถานการณ์ (Simulation) เช่น เกม Powerboat VR~\cite{powerboat2020} ซึ่งผู้ผลิตเกมระบุว่าผู้เล่นสามารถใช้คำสั่งเสียงในการควบคุมเรือแทนอุปกรณ์ควบคุมทั่วไป และเกม Tower! 3D~\cite{tower3d2017} ที่ผู้เล่นรับบทผู้ควบคุมการจราจรทางอากาศและใช้เสียงพูดตามมาตรฐานคำศัพท์การบินเพื่อสั่งการเครื่องบิน นอกจากเกมจำลองสถานการณ์แล้ว ยังมีเกมแนวผจญภัย เช่น In Verbis Virtus~\cite{inverbis2015} ที่ผู้เล่นร่ายเวทมนตร์ด้วยการพูดคาถาผ่านไมโครโฟน เกมเหล่านี้แสดงให้เห็นว่าการสั่งการด้วยเสียงสามารถนำมาใช้เป็นกลไกหลักของเกมได้ แต่ไม่ได้ออกแบบมาเพื่อการฝึกภาษาหรือลดความวิตกกังวลในการพูด

\subsection{สรุปเปรียบเทียบงานวิจัยและเกมที่เกี่ยวข้อง}

จากการทบทวนระบบและงานวิจัยข้างต้น สามารถสรุปเปรียบเทียบคุณสมบัติหลักได้ดังตารางที่ \ref{tab:compare_research}

\begin{table}[!htbp]
\centering
\caption{เปรียบเทียบคุณสมบัติของระบบและเกมที่เกี่ยวข้อง}
\label{tab:compare_research}
\begin{tabular}{|L{3.50cm}|L{2.00cm}|L{2.00cm}|L{2.00cm}|L{2.00cm}|}
\hline
\textbf{ระบบ/เกม} & \textbf{ควบคุมด้วยเสียง} & \textbf{NPC โต้ตอบด้วยเสียง} & \textbf{จำลองฟิสิกส์} & \textbf{ลดความกังวลภาษา} \\
\hline
Sadigzade~\cite{sadigzade2025} & \cmark & \cmark & -- & \cmark \\ \hline
Varma และคณะ~\cite{varma2025} & \cmark & -- & -- & -- \\ \hline
Zargham และคณะ~\cite{zargham2024} & \cmark & \cmark & -- & -- \\ \hline
Powerboat VR~\cite{powerboat2020} & \cmark & -- & \cmark & -- \\ \hline
Tower! 3D~\cite{tower3d2017} & \cmark & -- & -- & -- \\ \hline
In Verbis Virtus~\cite{inverbis2015} & \cmark & -- & -- & -- \\ \hline
\textbf{โครงงานนี้} & \cmark & \cmark & \cmark & ออกแบบเพื่อลด (รอประเมิน) \\
\hline
\end{tabular}
\end{table}

จากตารางที่ \ref{tab:compare_research} พบว่าระบบส่วนใหญ่แบ่งเป็น 2 กลุ่ม คือ เกมที่เน้นการเรียนรู้ภาษาแต่ไม่จำลองฟิสิกส์ และเกมจำลองสถานการณ์ที่ไม่ได้ออกแบบเพื่อลดความวิตกกังวลทางภาษา โครงงานนี้จึงออกแบบให้ครอบคลุมทุกด้าน

\subsection{เปรียบเทียบเทคโนโลยีเอนจินพัฒนาเกม (Game Engine)}

{\sloppy
ในการเลือกเทคโนโลยีสำหรับพัฒนาเกม ได้เปรียบเทียบ Game Engine ที่นิยม 3 ตัวเลือก ดังตารางที่ \ref{tab:compare_engine}
}

\begin{table}[!htbp]
\centering
\caption{เปรียบเทียบเทคโนโลยี Game Engine}
\label{tab:compare_engine}
\begin{tabular}{|L{2.96cm}|L{3.46cm}|L{3.46cm}|L{2.96cm}|}
\hline
\textbf{เกณฑ์} & \textbf{Unity} & \textbf{Unreal Engine} & \textbf{Godot} \\
\hline
ภาษาหลักที่ใช้ & C\# & C++ / Blueprint & GDScript / C\# \\ \hline
การรองรับ 3D & ดีมาก & ยอดเยี่ยม & ดี \\ \hline
การเชื่อมต่อ API ภายนอก & ง่าย มีไลบรารีรองรับมาก & ปานกลาง & ปานกลาง \\ \hline
แพ็กเกจ whisper.unity (whisper.cpp) สำหรับรู้จำเสียงภายในเครื่อง & ใช้ได้ทันที & ใช้ไม่ได้ (เป็นแพ็กเกจเฉพาะ Unity) & ใช้ไม่ได้ (เป็นแพ็กเกจเฉพาะ Unity) \\ \hline
ชุมชนผู้พัฒนา & ใหญ่มาก & ใหญ่มาก & กำลังเติบโต \\ \hline
\end{tabular}
\end{table}

จากตารางที่ \ref{tab:compare_engine} โครงงานนี้เลือกใช้ \textbf{Unity} (Unity 6) ร่วมกับภาษา \textbf{C\#} โดยปัจจัยหลักคือ Unity มีแพ็กเกจ whisper.unity ซึ่งเชื่อมต่อไลบรารี whisper.cpp ให้รันโมเดล Whisper ภายในเครื่องของผู้เล่นได้โดยตรง นอกจากนี้ Unity ยังมีไลบรารีสำหรับการจัดการเสียงที่ครอบคลุม และผู้จัดทำมีความคุ้นเคยในการเชื่อมต่อกับ API ภายนอก

% -----------------------------------------------------------
\section{หลักการและทฤษฎีที่เกี่ยวข้อง}
% -----------------------------------------------------------

\subsection{การเรียนรู้โดยใช้เกมเป็นฐาน (Digital Game-Based Learning: DGBL)}

เป็นการนำเกมดิจิทัลมาใช้เป็นเครื่องมือเพื่อส่งเสริมกระบวนการเรียนรู้ โดยอาศัยความสนุกสนานและความท้าทายเป็นตัวกระตุ้นแรงจูงใจภายใน (Intrinsic Motivation)~\cite{prensky2001} แนวคิดนี้เน้นให้ผู้เรียนเกิดการเรียนรู้เชิงประสบการณ์ (Experiential Learning) โครงงานนี้นำแนวคิดนี้มาเป็นแกนหลักในการเปลี่ยนรูปแบบการฝึกพูดภาษาอังกฤษให้กลายเป็นการเล่นเกมเพื่อบรรลุภารกิจ

\subsection{สมมติฐานตัวกรองทางอารมณ์ (Affective Filter Hypothesis)}

Krashen~\cite{krashen1982} อธิบายว่าตัวกรองทางอารมณ์ เช่น ความกลัว ความวิตกกังวล และการขาดความมั่นใจ เป็นอุปสรรคสำคัญที่ขัดขวางการเรียนรู้และการพูดภาษาที่สอง~\cite{sadigzade2025} โครงงานนี้นำสมมติฐานนี้มาประยุกต์ใช้โดยให้ตัวละครลูกเรือปัญญาประดิษฐ์ (AI Crew) เป็นผู้รับฟังคำสั่งแทนมนุษย์ ด้วยความคาดหวังว่าจะช่วยลดตัวกรองทางอารมณ์และลดความวิตกกังวลในการใช้ภาษาต่างประเทศ (Foreign Language Anxiety) ทำให้ผู้เล่นผ่อนคลายและกล้าเปล่งเสียงมากขึ้น ซึ่งจะประเมินในการทดลองกับกลุ่มตัวอย่างต่อไป

\subsection{ความวิตกกังวลในการใช้ภาษาต่างประเทศและแบบวัด FLCAS}

Horwitz, Horwitz และ Cope~\cite{horwitz1986} อธิบายว่าความวิตกกังวลในการใช้ภาษาต่างประเทศเป็นความวิตกกังวลเฉพาะสถานการณ์ที่เกิดขึ้นจากการเรียนและการใช้ภาษาต่างประเทศ ซึ่งแตกต่างจากความวิตกกังวลทั่วไป และส่งผลโดยตรงต่อการพูด Horwitz และคณะแบ่งความวิตกกังวลดังกล่าวออกเป็น 3 องค์ประกอบ ได้แก่ (1) ความประหม่าในการสื่อสาร (Communication Apprehension) (2) ความวิตกกังวลต่อการสอบ (Test Anxiety) และ (3) ความกลัวการถูกประเมินในทางลบ (Fear of Negative Evaluation) พร้อมทั้งพัฒนาแบบวัด Foreign Language Classroom Anxiety Scale (FLCAS) จำนวน 33 ข้อ แบบมาตราส่วนประมาณค่า 5 ระดับ มีคะแนนรวม 33--165 คะแนน โดยคะแนนสูงหมายถึงความวิตกกังวลสูง แบบวัดนี้เป็นเครื่องมือที่ใช้อย่างแพร่หลายในการศึกษาความวิตกกังวลทางภาษา

โครงงานนี้นำแนวคิดของ FLCAS มาใช้วัดผลว่าการฝึกพูดกับลูกเรือปัญญาประดิษฐ์ช่วยลดความวิตกกังวลในการพูดภาษาอังกฤษได้หรือไม่ โดยดัดแปลงให้เหมาะกับบริบทของเกม ดังนี้ (1) เลือกเฉพาะองค์ประกอบด้านความประหม่าในการสื่อสารและความกลัวการถูกประเมินในทางลบ ซึ่งเกี่ยวข้องกับการพูดโดยตรง โดยไม่รวมด้านความวิตกกังวลต่อการสอบ เนื่องจากเกมไม่มีการสอบ (2) เรียบเรียงข้อความใหม่ให้เป็นสถานการณ์การพูดภาษาอังกฤษทั่วไปแทนบริบทห้องเรียน และ (3) ลดจำนวนข้อเหลือ 12 ข้อ เพื่อให้ตอบซ้ำได้ทั้งก่อนและหลังเล่นเกมโดยไม่เป็นภาระแก่ผู้เข้าร่วม แบบวัดฉบับที่ใช้แสดงในภาคผนวก ก

\subsection{ทฤษฎีการเรียนรู้ตามสถานการณ์ (Situated Learning Theory)}

Lave และ Wenger~\cite{lave1991} เสนอว่าการเรียนรู้เกิดขึ้นได้ดีเมื่ออยู่ในบริบทหรือสถานการณ์ที่ความรู้นั้นนำไปใช้จริง โดยผู้เรียนค่อย ๆ เข้าร่วมในกิจกรรมของชุมชนผู้ปฏิบัติงาน โครงงานนี้นำแนวคิดดังกล่าวมาใช้โดยให้ผู้เล่นฝึกภาษาผ่านบริบทการสื่อสารทางทะเล (Maritime Communication) ได้แก่ การสั่งการทิศทาง การปรับความเร็ว และการแก้ปัญหาเฉพาะหน้า แทนการท่องจำวลีนอกบริบท

\subsection{วลีสื่อสารมาตรฐานทางทะเล (IMO Standard Marine Communication Phrases: SMCP)}

องค์การทางทะเลระหว่างประเทศ (International Maritime Organization: IMO) กำหนดวลีสื่อสารมาตรฐานทางทะเล (SMCP)~\cite{imo2002smcp} เพื่อให้การสื่อสารด้วยภาษาอังกฤษบนเรือและระหว่างเรือกับฝั่งมีความชัดเจน กระชับ และลดความกำกวม โดยเฉพาะในสถานการณ์ที่ผู้สื่อสารใช้ภาษาอังกฤษเป็นภาษาต่างประเทศ SMCP ครอบคลุมทั้งคำสั่งถือท้าย (Wheel Orders) คำสั่งเครื่องยนต์ (Engine Orders) และรูปแบบการทวนคำสั่งเพื่อยืนยันความเข้าใจ

คำสั่งในเกมของโครงงานนี้ยึดตามรูปแบบของ SMCP เช่น คำสั่งถือท้าย ``Starboard fifteen'' ``Midships'' และ ``hard-a-port'' โดยลูกเรือปัญญาประดิษฐ์ทวนคำสั่งด้วยรูปประโยคมาตรฐาน และเมื่อไม่เข้าใจคำสั่งจะตอบว่า ``Say again, Bridge. Over.'' ผู้เล่นจึงได้ฝึกใช้วลีที่ใช้จริงในงานเดินเรือแทนการสั่งการด้วยภาษาอิสระ

\subsection{เทคโนโลยีการรู้จำเสียง (Speech Recognition)}

การรู้จำเสียงคือกระบวนการแปลงคำสั่งเสียงพูดเป็นข้อความ (Speech-to-Text) โครงงานนี้เลือกใช้โมเดลรู้จำเสียง ได้แก่ Whisper~\cite{radford2023} เพื่อแปลงเสียงภาษาอังกฤษของผู้เล่นเป็นข้อความ โดยโครงงานนี้รันโมเดล ggml-base.en ภายในเครื่องของผู้เล่นผ่านไลบรารี whisper.cpp และประมวลผลทันทีหลังผู้เล่นปล่อยปุ่มวิทยุ จึงไม่ต้องส่งแฟ้มเสียงออกนอกเครื่อง ข้อความที่ได้จะส่งต่อให้ตัวแยกคำสั่ง (Command Parser) ภายในเกมเพื่อแยกคำสั่ง มีเพียงข้อความที่ถอดได้ในกรณีที่ตัวแยกคำสั่งไม่เข้าใจหรือผู้เล่นสนทนากับลูกเรือเท่านั้นที่ถูกส่งไปยัง OpenAI API

\subsection{แบบจำลองภาษาขนาดใหญ่สำหรับลูกเรือปัญญาประดิษฐ์ (GPT-4o mini)}

โครงงานนี้ใช้ OpenAI API กับโมเดล GPT-4o mini ใน 2 กรณี ได้แก่ (1) เป็นทางสำรองเมื่อตัวแยกคำสั่งภายในเกมไม่พบคำสั่งที่ชัดเจน โดยโมเดลจะตีความเจตนาของข้อความและส่งผลกลับเป็นคำสั่งที่เกมเข้าใจ และ (2) เมื่อผู้เล่นสนทนากับลูกเรือ โดยโมเดลจะสร้างข้อความตอบกลับของลูกเรือ ส่วนเสียงพูดของลูกเรือสร้างด้วยระบบสังเคราะห์เสียงของ Windows ภายในเครื่อง หากไม่สามารถใช้ API ได้ เช่น ไม่มีคีย์ API เกิดข้อผิดพลาด หรือหมดเวลารอคำตอบ ลูกเรือจะขอให้ผู้เล่นพูดซ้ำด้วยวลี ``Say again, Bridge. Over.''
